\subsection{Polynomial ambiguity in the invertible sector}\label{ssec:polyamb}
A weighted automaton $\mathcal P$ over $\Q$ consists of a finite state set $S$, an initial row
$\alpha\in\Q^{1\times S}$, a final column $\beta\in\Q^{S\times 1}$ and letter matrices
$A_a\in\Q^{S\times S}$, where $(A_a)_{pq}$ is the weight of the transition $p\xrightarrow{a}q$
($0$ if there is none); it computes $f(w)=\alpha A_w\beta$, $A_{a_1\cdots a_n}=A_{a_1}\cdots A_{a_n}$.
A \emph{path} labelled $w=a_1\cdots a_n$ from $p$ to $q$ is a sequence $p=s_0,\ldots,s_n=q$ with
$(A_{a_i})_{s_{i-1}s_i}\ne0$ for all $i$; it is \emph{accepting} if moreover $\alpha_{s_0}\ne0\ne\beta_{s_n}$.
The automaton is \emph{polynomially ambiguous} if the number of accepting paths labelled $w$ is
bounded by a polynomial in $|w|$~\cite{MPS26,PS24}. A series is polynomially ambiguous if some
polynomially ambiguous weighted automaton over $\Q$ computes it.

\begin{theorem}\label{thm:polyamb}
For an invertible rational series $f:\Sigma^*\to\Q$ the following are equivalent.
\begin{enumerate}
\item $f$ is polynomially ambiguous.
\item $\cG_f^\circ$ is triangularizable over $\Q$.
\item $\Gamma_f$ has a subgroup of finite index that is triangularizable over $\Q$.
\end{enumerate}
\end{theorem}

Together with Theorem~\ref{thm:main} and Corollary~\ref{cor:consequences} this settles the
relation between the two classes in the invertible sector.

\begin{corollary}\label{cor:polyamb}
An invertible series is CCRA-computable if and only if it is polynomially ambiguous and any two
unipotent elements of $\cG_f$ commute; equivalently, iff it is polynomially ambiguous and the
unipotent radical of $\cG_f^\circ$ is commutative. Given a rational weighted automaton, one can
decide in deterministic polynomial time whether its series is invertible and polynomially ambiguous.
\end{corollary}
\begin{proof}
The first claim follows from Theorem~\ref{thm:polyamb} and the equivalence (C1)$\Leftrightarrow$(C3) of
Theorem~\ref{thm:main}; if $\cG_f^\circ$ is triangularizable over $\Q$, every unipotent
element of $\cG_f$ lies in $\cG_f^\circ$ (A3), and the unipotent elements of $\cG_f^\circ$ form its
unipotent radical (Lemma~\ref{lem:split}). For the second claim, minimize and test invertibility
(Theorem~\ref{thm:decision}) and compute the block form of Section~\ref{ssec:block} (step~1 of the
algorithm there, in polynomial time); condition~(2) is condition~(a) of Appendix~\ref{sec:polytime}, which
is decidable in deterministic polynomial time by Lemma~\ref{lem:atest}.
\end{proof}

Thus in the invertible sector the CCRA series are exactly the polynomially ambiguous series for which
the unipotent radical of $\cG_f^\circ$ is commutative; the series of~\cite[Example~17]{MPS26} and of
Example~\ref{ex:heis} are polynomially ambiguous, with $\cG_f^\circ$ triangular over $\Q$ and
noncommutative unipotent radical. In particular, all eigenvalues of all word matrices of an
invertible polynomially ambiguous series are roots of rationals (elements of
$\Gamma_f\cap\cG_f^\circ$ are triangular over $\Q$, and $g^k\in\Gamma_f\cap\cG_f^\circ$ for
$g\in\Gamma_f$ and $k=[\Gamma_f:\Gamma_f\cap\cG_f^\circ]$), and $\cG_f^\circ$ modulo its unipotent radical
is a $\Q$-split torus (Lemma~\ref{lem:split}); a non-split torus, as for a rotation matrix of infinite order, never occurs.
This is in line with the one-letter statement of~\cite{BFLM22}.

\begin{remark}\label{rem:polyamb-ps}
Puch and Smertnig~\cite[Theorem~1.6]{PS24} characterize polynomial ambiguity of an invertible
series by a spectral condition on the word matrices (finitely generated spectrum and a
power-splitting property); with~\cite[Theorem~1.3, Proposition~9.4]{PS24} this is equivalent to the
representation being an epimorphic image of a block-triangular representation with monomial diagonal
blocks. That structural form is closely related to condition~(3) of Theorem~\ref{thm:polyamb}, since a
block-triangular group with monomial diagonal blocks has a triangularizable subgroup of finite index.
The sufficiency direction there uses unit equations, and decidability is obtained through
Zariski-closure computations. Theorem~\ref{thm:polyamb} gives an intrinsic condition on the Zariski
closure of the minimal group, its proof below is direct (it uses neither their result nor unit
equations), and it yields the polynomial-time bound of Corollary~\ref{cor:polyamb}. The proof of Theorem~\ref{thm:polyamb}
uses only that $\Q$ has characteristic zero and holds verbatim over any field of characteristic zero,
with automata and triangularizability taken over that field; the complexity claim is specific to $\Q$.
\end{remark}

\begin{proof}[Proof of Theorem~\ref{thm:polyamb}: (2)$\Leftrightarrow$(3) and (2)$\Rightarrow$(1)]
Write $\Gamma=\Gamma_f$, $\cG=\cG_f$ and $\Gamma^\circ=\Gamma\cap\cG^\circ$, a normal subgroup of
finite index (Lemma~\ref{lem:monoid}). If (2) holds, $\Gamma^\circ$ fixes the rational flag fixed by
$\cG^\circ$. If $\Gamma_1$ is as in (3), its closure has finite index in $\cG$, hence contains
$\cG^\circ$ (A1), and fixes the rational flag fixed by $\Gamma_1$.

Assume (2) and choose $h_0\in\GL_d(\Q)$ such that $h_0^{-1}gh_0$ is upper triangular for all
$g\in\cG^\circ$. With $F=\Gamma/\Gamma^\circ$, $\varphi$, the words $r_g$ and the cocycle
$c(g,a)=M_{r_g}M_aM_{r_{g\varphi(a)}}^{-1}\in\Gamma^\circ$ as in Section~\ref{sec:normalform} (their
construction and~\eqref{eq:cocycle} use only Lemma~\ref{lem:monoid}, not condition~(C3)), the
matrices $c'(g,a)=h_0^{-1}c(g,a)h_0$ are rational and upper triangular, and by~\eqref{eq:cocycle}
\[
f(w)=(\lambda h_0)\,c'(g_0,a_1)\cdots c'(g_{n-1},a_n)\,(h_0^{-1}M_{r_{g_n}}\gamma).
\]
Let $\mathcal P$ have states $F\times\{1,\ldots,d\}$, a transition $(g,i)\xrightarrow{a}(g\varphi(a),j)$
of weight $c'(g,a)_{ij}$ whenever this entry is nonzero, initial weights $(\lambda h_0)_i$ on $(1,i)$ and
final weights $(h_0^{-1}M_{r_g}\gamma)_j$ on $(g,j)$. Its letter matrices are block permutation matrices
with blocks $c'(g,a)$, so $\mathcal P$ computes $f$ by the displayed formula. On a word of length $n$
the first components along an accepting path are $g_0,\ldots,g_n$, and the second components form a
nondecreasing sequence $i_0\le\cdots\le i_n$ in $\{1,\ldots,d\}$, because each $c'(g,a)$ is upper
triangular. Such a sequence is determined by the number of occurrences of each index, so there are at most $(n+2)^d$ accepting paths, and $\mathcal P$ is
polynomially ambiguous.
\end{proof}

For (1)$\Rightarrow$(2), fix a polynomially ambiguous automaton $\mathcal P$ computing $f$, which we may
assume \emph{trim}: every state lies on an accepting path (removing the other states changes neither
$f$ nor the accepting paths).

\begin{lemma}\label{lem:uniquepath}
Let $C$ be a strongly connected component of the transition graph of $\mathcal P$. For every word
$w$ and all $p,q\in C$ there is at most one path labelled $w$ from $p$ to $q$, and its states lie in
$C$.
\end{lemma}
\begin{proof}
Every state of such a path is reachable from $p$ and reaches $q$, so it lies in $C$. Suppose that
$\pi_1\ne\pi_2$ are two paths labelled $w$ from $p$ to $q$; then $|w|\ge1$. Let $\tau$ be a path from
$q$ to $p$, labelled $v$ (empty if $p=q$). Then $\pi_1\tau$ and $\pi_2\tau$ are distinct paths labelled $wv$ from $p$ to $p$. As
$\mathcal P$ is trim, there are paths labelled $u$ from a state with nonzero initial weight to $p$ and
labelled $u'$ from $p$ to a state with nonzero final weight. Choosing one of the two cycles for each
of the $k$ blocks gives $2^k$ accepting paths labelled $u(wv)^ku'$, pairwise distinct because the blocks
have the same length and are aligned; as $|u(wv)^ku'|$ grows linearly in $k$, this contradicts polynomial
ambiguity.
\end{proof}

\begin{lemma}[Component lemma]\label{lem:component}
Let $C$ be a finite set and $\rho(a)\in\Q^{C\times C}$ for $a\in\Sigma$, with
$\rho(a_1\cdots a_n)=\rho(a_1)\cdots\rho(a_n)$. Define paths as above and assume:
\begin{itemize}
\item[(U)] for every word $w$ and all $p,q\in C$ there is at most one path labelled $w$ from $p$ to $q$.
\end{itemize}
Let $Y\subseteq X\subseteq\Q^C$ be subspaces invariant under every $\rho(a)$ (acting on columns), such
that every $\rho(a)$ induces an invertible operator $\sigma(a)$ on $D=X/Y$, and let $\cH$ be the
Zariski closure of $\{\sigma(w):w\in\Sigma^*\}$ in $\GL(D)$. Then $\cH^\circ$ is a torus split over
$\Q$; in particular $\cH^\circ$ is diagonalizable over $\Q$.
\end{lemma}
\begin{proof}
\emph{Supports.} By (U), $\rho(w)_{pq}$ is the product of the weights along the unique path labelled
$w$ from $p$ to $q$ if it exists, and $0$ otherwise. Let $B_w\subseteq C\times C$ be the support of
$\rho(w)$. Then $B_{uv}=B_u\circ B_v$ (composition of relations), so $w\mapsto B_w$ is a morphism
onto a finite monoid $\mathfrak B$ of relations. For words $z_1,\ldots,z_k$ the sequences
$p=s_0,\ldots,s_k=q$ with $(s_{i-1},s_i)\in B_{z_i}$ for all $i$ correspond injectively to paths
labelled $z_1\cdots z_k$ from $p$ to $q$, so by (U) there is at most one such sequence, and
\begin{equation}\label{eq:singleterm}
\rho(z_1\cdots z_k)_{pq}=\textstyle\prod_{i=1}^k\rho(z_i)_{s_{i-1}s_i}
\end{equation}
for that sequence if it exists, and $0$ otherwise.

\emph{An idempotent with a group around it.} The intersection $\mathfrak J$ of all two-sided ideals of
$\mathfrak B$ is nonempty (it contains the product of all of them), hence the minimum ideal, and
$\mathfrak B^1x\mathfrak B^1=\mathfrak J$ for every $x\in\mathfrak J$. Some power of an element of
$\mathfrak J$ is an idempotent $E\in\mathfrak J$; write $E=B_e$ for a word $e$. The finite monoid
$\mathfrak M=E\mathfrak BE$ with identity $E$ is a group: for $x\in\mathfrak M$ we have
$E=sxt$ with $s,t\in\mathfrak B^1$, hence $E=(EsE)x(EtE)$; so $t'(s'x)=E$ for
$s'=EsE$, $t'=EtE\in\mathfrak M$, because in a finite monoid $yz=1$ implies $zy=1$; thus $x$ has a
left inverse, hence is a unit. Let $n=|\mathfrak M|$. Then $B_{(exe)^n}=(EB_xE)^n=E$ for every word
$x$. Put $L=\{z\in\Sigma^*:B_z=E\}$, a semigroup containing $e$ and all $(exe)^n$. If $D\ne0$ then
$\sigma(e)$ is invertible, so $\rho(e)\ne0$ and $E\ne\emptyset$; if $D=0$ there is nothing to prove.

\emph{Fixed points of $E$.} Let $\Theta=\{\theta\in C:(\theta,\theta)\in E\}$. By the uniqueness statement above with
all $B_{z_i}=E$, for every $k\ge1$ and all $p,q$ there is at most one sequence of length $k$ from $p$
to $q$ with consecutive pairs in $E$ (an \emph{$E$-sequence}). Hence:
(i) if $\theta\ne\theta'$ are in $\Theta$ then $(\theta,\theta')\notin E$, since otherwise $\theta,\theta,\theta'$ and $\theta,\theta',\theta'$ are two
$E$-sequences;
(ii) for $(p,q)\in E$ there is exactly one $\theta\in\Theta$ with $(p,\theta),(\theta,q)\in E$. Uniqueness is clear.
For existence, $E=E^{k}$ for $k=|C|+1$ yields an $E$-sequence $s_0=p,\ldots,s_k=q$, in which some
state repeats, $s_i=s_j$ with $i<j$; then $(s_i,s_i)\in E^{j-i}=E$, so $\theta=s_i\in\Theta$, and
$(p,\theta)\in E$ (if $i=0$ then $p=\theta$) and $(\theta,q)\in E$ (if $j=k$ then $\theta=q$).

\emph{Factorization through a diagonal character.} Let $\iota\in\Q^{C\times\Theta}$ be the inclusion of
coordinates. For $z\in L$ the matrix $\chi(z)=\iota^{\mathsf T}\rho(z)\iota$ is diagonal by (i), with
nonzero diagonal entries. For $z,z'\in L$, \eqref{eq:singleterm} and the uniqueness of
$E$-sequences give $\chi(zz')=\chi(z)\chi(z')$: the only $E$-sequence $\theta,r,\theta$ is $r=\theta$. Moreover, for
$z,y,z'\in L$,
\begin{equation}\label{eq:factor}
\rho(zyz')=\rho(z)\,\iota\,\chi(y)\,\iota^{\mathsf T}\rho(z').
\end{equation}
Indeed, the $(p,q)$ entry of the right side is $\sum_{\theta\in\Theta}\rho(z)_{p\theta}\chi(y)_{\theta\theta}\rho(z')_{\theta q}$,
whose nonzero terms correspond to $E$-sequences $p,\theta,\theta,q$; the left side is given by
\eqref{eq:singleterm} with $E$-sequences $p,r,s,q$ of length three. If $(p,q)\notin E=E^3$, both sides
vanish. If $(p,q)\in E$, the unique $\theta$ of (ii) gives an $E$-sequence $p,\theta,\theta,q$, which is therefore
the only $E$-sequence of length three from $p$ to $q$, and both sides equal
$\rho(z)_{p\theta}\rho(y)_{\theta\theta}\rho(z')_{\theta q}$.

\emph{A morphism from a $\Q$-split diagonalizable group onto the closure.} Let $\mathcal D\subseteq\Q^{\Theta\times\Theta}$ be
the linear span of $\chi(L)$, a space of diagonal matrices. By~\eqref{eq:factor} with $z=z'=e$, the
linear map $\delta\mapsto\rho(e)\iota\delta\iota^{\mathsf T}\rho(e)$ sends $\chi(y)$ to $\rho(eye)$ for
$y\in L$; hence on $\mathcal D$ it takes values preserving $X$ and $Y$ and induces a linear map
$\Lambda:\mathcal D\to\operatorname{End}(D)$ with $\Lambda(\chi(y))=\sigma(eye)=\sigma(e)\sigma(y)\sigma(e)$.
The linear map $\Psi(\delta)=\sigma(e)^{-1}\Lambda(\delta)\sigma(e)^{-1}$, defined over $\Q$ and
extended to $\mathcal D\otimes\Qb$, satisfies $\Psi(\chi(y))=\sigma(y)$ for $y\in L$. Let $\mathcal T$
be the Zariski closure of $\chi(L)$ in the diagonal torus of $\GL_{\Theta}$; it lies in
$\mathcal D\otimes\Qb$ and is a closed subgroup by the proof of Lemma~\ref{lem:monoid}, which applies to
subsemigroups of $\GL$ with one more line: from $\bar S=x\bar S$ and $x\in\bar S$ we get $x=xy$ with
$y\in\bar S$, so $1=y\in\bar S$ and then $x^{-1}\in\bar S$. The polynomial identity $\Psi(ab)=\Psi(a)\Psi(b)$ holds on
$\chi(L)\times\chi(L)$, hence on its closure $\mathcal T\times\mathcal T$; and $\Psi(1)=1$, because
$\Psi(1)\Psi(\chi(e))=\Psi(\chi(e))=\sigma(e)$ is invertible; then $\Psi(a)\Psi(a^{-1})=1$. So $\Psi:\mathcal T\to\GL(D)$ is a
morphism of algebraic groups defined over $\Q$. Its image is closed (A2), contains $\sigma(L)$, and
lies in the closure $\cK$ of $\sigma(L)$; hence $\cK=\Psi(\mathcal T)$ and $\cK^\circ=\Psi(\mathcal T^\circ)$.
The torus $\mathcal T^\circ$ is a subtorus of a diagonal torus defined over $\Q$, hence split over $\Q$,
and so is its image $\cK^\circ$ under the $\Q$-morphism $\Psi$ (characters of $\cK^\circ$ pull back
injectively and Galois-equivariantly to characters of $\mathcal T^\circ$).

\emph{The closure contains $\cH^\circ$.} This is the argument of Lemma~\ref{lem:dense}. The group $\cH$ is
the closure of $\sigma(\Sigma^*)$ (Lemma~\ref{lem:monoid}), and $\cK\subseteq\cH$ is a closed subgroup.
Let $g=\sigma(e)$ and $\Xi(h)=(ghg)^n$. For every word $x$, $\Xi(\sigma(x))=\sigma((exe)^n)\in\cK$,
so $\Xi(\cH)\subseteq\cK$ by density; for $k\in\cH^\circ$, $\Xi(g^{-1}kg^{-1})=k^n\in\cK$. The closure
of the group generated by all $k^n$, $k\in\cH^\circ$, is normal in $\cH^\circ$ and contained in $\cK$,
and the quotient of $\cH^\circ$ by it is connected of finite exponent, hence trivial (A4). Thus
$\cH^\circ\subseteq\cK\subseteq\cH$, so $\cH^\circ=\cK^\circ$ is a $\Q$-split torus, and a $\Q$-split torus
acting on a rational vector space is diagonalizable over $\Q$~\cite[\S8]{Bor91}.
\end{proof}

\begin{proof}[Proof of Theorem~\ref{thm:polyamb}: (1)$\Rightarrow$(2)]
Let $\mathcal P=(\alpha,(A_a)_a,\beta)$ be trim and polynomially ambiguous, computing $f$. Number its
strongly connected components $C_1,\ldots,C_r$ so that every transition goes from $C_i$ to $C_j$ with
$i\le j$. Then $V_j=\spn\{e_q:q\in C_1\cup\cdots\cup C_j\}$ is invariant under every $A_a$ acting on
columns, and $A_w$ induces on $V_j/V_{j-1}\cong\Q^{C_j}$ the diagonal block
$\rho_j(w)=(A_w)_{C_jC_j}$; the $\rho_j$ satisfy (U) by Lemma~\ref{lem:uniquepath}.

By Section~\ref{sec:prelim}, with $W=\spn\{A_w\beta\}$ and $N=\{h\in W:\alpha A_vh=0\ \forall v\}$, the
operators induced on $W/N$ form a representation isomorphic to the minimal one; identify $W/N$ with
$\Q^d$. Put $W_j=W\cap V_j$ and $\Phi_j=(W_j+N)/N$, so $0=\Phi_0\subseteq\cdots\subseteq \Phi_r=\Q^d$ are
rational subspaces invariant under every $M_a$, hence under $\Gamma_f$ and $\cG_f$. Restricting
coordinates to $C_j$ maps $W_j$ equivariantly onto $X_j\subseteq\Q^{C_j}$ with kernel $W_{j-1}$, and
the surjection $W_j\to \Phi_j/\Phi_{j-1}$ has kernel $W_j\cap(W_{j-1}+N)\supseteq W_{j-1}$. Hence
$\Phi_j/\Phi_{j-1}\cong X_j/Y_j$, where $Y_j$ is the image of $W_j\cap(W_{j-1}+N)$, and this isomorphism
intertwines the operator induced by $M_w$ with the operator $\sigma_j(w)$ induced by $\rho_j(w)$. As
$M_w$ is invertible, so is $\sigma_j(w)$. By Lemma~\ref{lem:component} the closure $\cH_j$ of
$\{\sigma_j(w)\}$ has an identity component diagonalizable over $\Q$. The induced-action morphism
$\varpi_j$ from $\cG_f$ to $\GL(\Phi_j/\Phi_{j-1})$ maps $\cG_f$ onto $\cH_j$ and $\cG_f^\circ$ onto $\cH_j^\circ$
(A2). Choosing in each $\Phi_j/\Phi_{j-1}$ a rational basis diagonalizing $\cH_j^\circ$ and lifting, we obtain
a complete rational flag of $\Q^d$ fixed by $\cG_f^\circ$.
\end{proof}

The proof shows more: $\cG_f^\circ$ acts on each graded piece $\Phi_j/\Phi_{j-1}$ through a $\Q$-split torus,
so the unipotent radical of $\cG_f^\circ$ is exactly the subgroup acting trivially on all of them
(that subgroup is normal and consists of unipotent elements, and unipotent elements act trivially
through a torus).
Polynomial ambiguity thus controls the reductive part of $\cG_f^\circ$, whereas copylessness
additionally forces the unipotent radical to be commutative.
